\scanpage{385}
\mspage{301}
% archive sheet number: 384

cofinale. Donc il reste à \deleted{définir} construire
$i_\alpha$ pour $\alpha$ ordinal limite. Or
l'ens des $i_\beta$ pour $\beta<\alpha$ est de
cardinal $<\pi$, et comme \inserted{par a)} $I$ est g\textsuperscript{r}de
devant $\pi$ (il suffit qu'il le soit devant
tout cardinal $<\pi$) il existe un
$i\in\operatorname{Ob}I$ qui \unsure{majore} tous les $i_\beta$.
Choisissons des
\[
i_\beta\xrightarrow{\unsure{w|u}_\beta}i,
\]
On n'aura pas a priori comm. des

\diagram{\begin{tikzcd}[ampersand replacement=\&]
\deleted{\unsure{v_\alpha\tau_\beta}}\phantom{i}\\[-1ex]
i_{\beta'}\arrow[r,"{\tau_{\beta,\beta'}}"]\arrow[rr,bend right=17,"{u_{\beta'}}"'] \& i_\beta\arrow[r,"{u_\beta}"] \& i\mathpunct{,}
\end{tikzcd}}

mais comme l'ens des doubles-flèches

\diagram{\begin{tikzcd}[ampersand replacement=\&]
i_{\beta'}\arrow[r,shift left=0.5ex,"{u_\beta\tau_{\beta,\beta'}}"]\arrow[r,dash,shift right=0.5ex,"{u_{\beta'}}"'] \& i
\end{tikzcd}}

est de cardinal $<\pi$, $I$ étant g\textsuperscript{r}de devant $\pi$,
il existe $i\longrightarrow i'$ qui les égalise
tous. Comme $J_{\unsure{\varphi|\phi}(\alpha)}$ est cofinal dans
$I$, il existe $i''\in J_{\unsure{\varphi\alpha|\varphi(\alpha)}}$, et $i'\unsure{\text{---}|\longrightarrow}i''$.
On prend $i_\alpha=i''$, avec les hom composés

\diagram{\begin{tikzcd}[ampersand replacement=\&]
i_\beta \arrow[r,"{u_\beta}"]\arrow[rrr,bend right=12,"{\tau_{\alpha,\beta}}"'] \& \unsure{i_\beta|u_\beta|i}\arrow[r] \& i'\arrow[r,dash] \& i^{\unsure{\prime\prime|\prime}}
\end{tikzcd}}
